%%%%%%%%%%%%%%%%%%%%%%%%%%%%%%%%%%%%%%%%%%%%%%%%
%
% Strath PhD Thesis Template
%  by Jethro Browell [jethro.browell@strath.ac.uk]
%
%  Guidelines for thesis format, submission and content are found in
%  General and Course Regulations for Graduate and Postgraduate
%  Awards and Degrees, section 20.6.
%
%  Using .eps or .pdf is recomended to prduce high quality figures etc.
%
%  The Strathclyde logo can be found in other formats at www.strath.ac.uk.
%
%%%%%%%%%%%%%%%%%%%%%%%%%%%%%%%%%%%%%%%%%%%%%%%%

\documentclass[a4paper,oneside,11pt]{book}

\usepackage{amsbsy}
\usepackage{amsmath}
\usepackage{amsfonts}
\usepackage{graphicx}
\usepackage{multirow}
\usepackage{mathrsfs}
\usepackage{xcolor}
\usepackage[hidelinks]{hyperref}
\usepackage{cite}
\usepackage{enumitem}
\usepackage{epsfig}
\usepackage{caption}
\usepackage{subcaption}
\usepackage[strict]{changepage}

\usepackage{amssymb}
\usepackage{catchfilebetweentags}
\usepackage{cmll}
\usepackage[utf8]{inputenc}
\usepackage{mathtools}
\usepackage{mathpartir}
\usepackage{newunicodechar}
\usepackage{stmaryrd}
\usepackage{todonotes}

\usepackage[conor]{agda}

\newunicodechar{λ}{\ensuremath{\mathnormal\lambda}}
\newunicodechar{→}{\ensuremath{\mathnormal\to}}
\newunicodechar{∀}{\ensuremath{\mathnormal\forall}}
\newunicodechar{ι}{\ensuremath{\mathnormal\iota}}
\newunicodechar{·}{\ensuremath{\mathnormal\cdot}}
\newunicodechar{⊸}{\ensuremath{\mathnormal\multimap}}
\newunicodechar{⊕}{\ensuremath{\mathnormal\oplus}}
\newunicodechar{─}{\textrm{---}}
\newunicodechar{ᶜ}{\ensuremath{\mathnormal{^c}}}
\newunicodechar{⊢}{\ensuremath{\mathnormal\vdash}}
\newunicodechar{∧}{\ensuremath{\mathnormal\wedge}}
\newunicodechar{⟦}{\ensuremath{\mathnormal\llbracket}}
\newunicodechar{⟧}{\ensuremath{\mathnormal\rrbracket}}
\newunicodechar{∩}{\ensuremath{\mathnormal\cap}}
\newunicodechar{✴}{\ensuremath{\mathnormal*}}
\newunicodechar{ℓ}{\ensuremath{\mathnormal\ell}}
\newunicodechar{Γ}{\ensuremath{\mathnormal\Gamma}}
\newunicodechar{Δ}{\ensuremath{\mathnormal\Delta}}
\newunicodechar{Σ}{\ensuremath{\mathnormal\Sigma}}
\newunicodechar{∈}{\ensuremath{\mathnormal\in}}
\newunicodechar{′}{\ensuremath{\mathnormal'}}
\newunicodechar{≡}{\ensuremath{\mathnormal\equiv}}
\newunicodechar{⊤}{\ensuremath{\mathnormal\top}}
\newunicodechar{⊥}{\ensuremath{\mathnormal\bot}}
\newunicodechar{⇒}{\ensuremath{\mathnormal\Rightarrow}}
\newunicodechar{▹}{\ensuremath{\mathnormal\triangleright}}
\newunicodechar{₁}{\ensuremath{\mathnormal{_1}}}
\newunicodechar{ℑ}{\ensuremath{\mathfrak I}}
\newunicodechar{□}{\ensuremath{\mathnormal\Box}}

\definecolor{use}{HTML}{008000}
\newcommand\gr[1]{{\color{use}#1}}
\newcommand\grctx[1]{\gr{\mathcal{#1}}}
\newcommand\grP{\grctx P}
\newcommand\grQ{\grctx Q}
\newcommand\grR{\grctx R}
\newcommand\grPprime{\grP\gr'}
\newcommand\grQprime{\grQ\gr'}
\newcommand\grRprime{\grR\gr'}
\newcommand\name{\ensuremath{\lambda\grR}}
\newcommand\grctxsub[2]{\grctx{#1}_{\gr{#2}}}
\newcommand\grPe{\grctxsub P e}
\newcommand\grPf{\grctxsub P f}
\newcommand\grPl{\grctxsub P l}
\newcommand\grPr{\grctxsub P r}
\newcommand\grQe{\grctxsub Q e}
\newcommand\grQf{\grctxsub Q f}
\newcommand\grQl{\grctxsub Q l}
\newcommand\grQr{\grctxsub Q r}
\newcommand\grRl{\grctxsub R l}
\newcommand\grRr{\grctxsub R r}
\newcommand\ps{\mathit{ps}}
\newcommand\qs{\mathit{qs}}
\newcommand\rs{\mathit{rs}}
\newcommand\dotto{\mathrel{\dot\to}}
\newcommand\dotlr{\mathrel{\dot\leftrightarrow}}
\newcommand\dottimes{\mathbin{\dot\times}}
%\newcommand\dotplus{\mathbin{\dot+}}
\newcommand\wand{\mathrel{\mathord{-}\hspace{-0.75ex}*}}
\newcommand\sep{\mathbin{*}}
\newcommand\Boxzp{\Box^{0{+}}}
\newcommand\Boxzpt{\Box^{0{+}{*}}}
\providecommand\U{}
\renewcommand\U{\mathcal U}
\newcommand\V{\mathcal V}
\newcommand\W{\mathcal W}
\providecommand\C{}
\renewcommand\C{\mathcal C}
\newcommand\sqin{\mathrel{\mathrlap{\sqsubset}{\mathord{-}}}}
\newcommand\sqni{\mathrel{\mathrlap{\sqsupset}{\mathord{-}}}}
\newcommand\Ann{\mathscr R}

\renewcommand\land{~\wedge~}
\renewcommand\lor{~\vee~}
\newcommand\rel{\mathrel{\mathord{\to}\hspace{-2.25ex}+}}

\DeclareMathOperator\Set{Set}
\DeclareMathOperator\obj{Obj}
\DeclareMathOperator\List{List}
\let\hom\relax
\DeclareMathOperator\hom{Hom}
\DeclareMathOperator\id{id}
\DeclareMathOperator\sub{Sub}

\newcommand\env[1]{\stackrel{#1}\Longrightarrow}

\usepackage{mleftright}
\newcommand\lr[3]{\mleft#1{#2}\mright#3}
\newcommand\sem[1]{\lr\llbracket{#1}\rrbracket}
\newcommand\size[1]{\lr\lvert{#1}\rvert}
\newcommand\plr[1]{\lr({#1})}
\newcommand\blr[1]{\lr[{#1}]}
\newcommand\forallb[1]{\forall\blr{~#1~}}
\newcommand\alr[1]{\lr\langle{#1}\rangle}
\newcommand\bra[1]{\lr\langle{#1}\rvert}
\newcommand\ket[1]{\lr\lvert{#1}\rangle}

\newcommand\leO{\;{\leq}0}
\newcommand\leI{\;{\leq}1}

\DeclareMathOperator{\lin}{lin}
\DeclareMathOperator{\intu}{int}
\newcommand\lnl{\ensuremath{\mathrm{L/nL}}}
\newcommand\vdashL{\mathrel{\vdash_{\mathcal L}}}
\newcommand\vdashC{\mathrel{\vdash_{\mathcal C}}}

\ebproofnewstyle{comb}{separation=0.75em}
\ebproofset{right label template=\TirName{\inserttext}}

\newenvironment{eqns}{\begin{array}{r@{\hspace{0.3em}}c@{\hspace{0.3em}}l}}{\end{array}}

\newcommand\PsiDot[1]{%
  \AgdaBound{$\Psi$}\AgdaSpace{}\AgdaSymbol{.}\AgdaField{#1}}

\newcommand\qto{\mathbin{`\!\to}}

\newcommand\Lock{\text{\faLock}}

\newcommand\oiw{\mathrm{\gr{01\upomega}}}

\newcommand\colour{%
{\color{AgdaField}co}{\color{AgdaFunction}lo}{\color{AgdaDatatype}ur}%
}


% Page Margins - Strath Requirement
\usepackage[left=4cm,right=2.5cm,top=2cm,bottom=4cm,includehead,includefoot,headheight=15pt]{geometry}

% Page Headers
\usepackage{fancyhdr}
\fancyhf{}
\renewcommand{\headrulewidth}{0pt} % optional
%\fancyhead[L]{\nouppercase{\leftmark} \hfill Section \nouppercase{\rightmark}}
\fancyhead[L]{\nouppercase{\leftmark}}
\cfoot{\thepage}
\pagestyle{fancy}

% Draft Watermark
%\usepackage[draft=true,allpages=true,fontfamily=cmr,angle=90,scale=0.1,mark={\fboxsep=35pt\fboxrule=0pt\relax\fbox{-- DRAFT -- \today~--}},xcoord=-80,ycoord=-20]{draftwatermark}


% Line Spacing
%\def\baselinestretch{1.5}
\usepackage{setspace}
\setstretch{1.5}


% Place UoS Logo on Title Page (this package modifies the "\maketitel" command.)
\usepackage{titling}
\pretitle{%
  \begin{flushright}
  \vspace{-9.5cm}
%  \includegraphics[width=5cm,natwidth=472,natheight=531]{logo} \\[7cm]
  \includegraphics[width=5cm]{logo} \\[6cm]
  \end{flushright}
  \begin{center}
  \LARGE
}
\posttitle{\end{center}}

\title{PhD Title \\ PhD Thesis}
\author{James Wood
\\ \small Mathematically Structured Programming\\[-0.8ex]
\small Computer and Information Sciences\\[-0.8ex]
\small University of Strathclyde, Glasgow\\
}




%%%%%%%%%%%%%%%%%%%%%%%%%%%%%%%%%%%%%%%%%%%%%%%%%%%%%%%%%%%%%%
\begin{document}

\maketitle


%%%%%%%%%%%%%%%%%%%%%%%%%%%%%%%%%%%%%%%%%%%%%%%%%%%%%%%%%%%%%%
\frontmatter
%%%%%%%%%%%%%%%%%%%%%%%%%%%%%%%%%%%%%%%%%%%%%%%%%%%%%%%%%%%%%%

% Declaration
\topskip0pt
\vspace*{\fill}
\noindent
\begin{quote}
  \centering
  This thesis is the result of the author's original research. It has been composed by the author and has not been previously submitted for examination which has led to the award of a degree. \\[5pt]
  %
  The copyright of this thesis belongs to the author under the terms of the United Kingdom Copyright Acts as qualified by University of Strathclyde Regulation 3.50. Due acknowledgement must always be made of the use of any material contained in, or derived from, this thesis. \\[5pt]
  %
  %Signed: \\
  %Date:
\end{quote}
\vspace*{\fill}



\chapter{Abstract}



\tableofcontents

\addcontentsline{toc}{chapter}{List of Figures}
\listoffigures

\addcontentsline{toc}{chapter}{List of Tables}
\listoftables



\chapter{Preface/Acknowledgements}
I would like to acknowledge...

This document is adapted from the template by Jethro Browell
(\url{https://www.overleaf.com/latex/templates/thesis-template-for-university-of-strathclyde/nfnrnmjqyxqg}),
which was licensed under CC BY 4.0.



%%%%%%%%%%%%%%%%%%%%%%%%%%%%%%%%%%%%%%%%%%%%%%%%%%%%%%%%%%%%%%
\mainmatter
%%%%%%%%%%%%%%%%%%%%%%%%%%%%%%%%%%%%%%%%%%%%%%%%%%%%%%%%%%%%%%


\chapter{Introduction}

\chapter{Simply typed $\lambda$-calculus}
  \section{Basics}
  \section{Kits}
  \section{Natural deduction vs sequent calculus}
  In a seminal paper~\cite{Gentzen64}, Gerhard Gentzen introduces two syntactic
paradigms which remain among the most studied to this day.
These paradigms are \emph{natural deduction} and \emph{sequent calculus}, as
exemplified by natural deduction calculi NJ and NK, and sequent calculi LJ and
LK.
Restricting attention to just the intuitionistic systems NJ and LJ, these
actually differ in two orthogonal ways, which I shall prise apart in this
section.
The simpler distinction is that the logical rules in NJ are introduction and
elimination rules, whereas the logical rules in LJ are left and right rules.
But the more important distinction for this thesis is that where NJ has
assumptions, LJ has sequents explictly manipulated by structural rules.
I take the latter distinction to define natural deduction and sequent calculus,
and wherever I need to make the former distinction I shall speak of
\emph{intro-elim systems} and \emph{left-right systems}.

I will use the rest of this section as follows.
First, I introduce NJ and LJ, and some of their basic metatheory.
Then, I will give examples of systems intermediate between Gentzen's natural
deduction and sequent calculi: a left-right natural deduction calculus
$\mu\tilde\mu$, and an intro-elim sequent calculus BBdPH\@.
Both of these examples will reappear in later chapters.

\subsection{Intro-elim natural deduction: NJ}

\subsection{Left-right sequent calculus: LJ}

\subsection{Left-right natural deduction: $\mu\tilde\mu$}
The $\mu\tilde\mu$-calculus~\cite{CH00} (also known as
$\overline\lambda\mu\tilde\mu$ or system L, and closely related to Wadler's
dual calculus~\cite{Wadler03}) can be seen as an adaptation of natural deduction
to classical logic.
Though originally presented as a sequent calculus, the underlying natural
deduction calculus was later given by Herbelin~\cite[p.\ 12]{herbelin-hab}, and
I will follow the latter.
While Gentzen gave a natural deduction calculus NK for classical logic, NK
relies on adding the \emph{axiom} of excluded middle.
As axioms are not systematic parts of the calculus, they can hinder or break
metatheoretic properties like normalisation.
In contrast, the $\mu\tilde\mu$-calculus allows us to \emph{derive} excluded
middle from entirely systematic components.

In NJ, a derivation of $A$ from assumptions $\Gamma$ tells us that if each
formula in $\Gamma$ is \emph{true}, then $A$ is also \emph{true}.
The $\mu\tilde\mu$-calculus generalises the picture by allowing us to have
both \emph{true} and \emph{false} assumptions, and allowing us to conclude that
some $A$ is \emph{true}, that some $A$ is \emph{false}, or that we have reached
a contradiction.
% A similar judgement of contradiction appears in Prawitz' classical natural
% deduction calculus~\cite{Prawitz65}.
Following Herbelin, we notate the judgement that $A$ is true as ${}\vdash A$,
that $A$ is false as $A \vdash{}$, and of contradiction as $\vdash$.
The only way to derive a contradiction is to find some $A$ such that
${}\vdash A$ and $A \vdash{}$.
Meanwhile, we can derive ${}\vdash A$ by assuming $A \vdash{}$ and deriving
$\vdash$, i.e., we can prove $A$ by assuming that $A$ is false and deriving a
contradiction.
Dually, we can derive $A \vdash{}$ by assuming ${}\vdash A$ and deriving
$\vdash$.
These three methods of derivation are encoded in the following rules.

\begin{mathpar}
  \inferrule*[right=Cut]
  {{}\vdash A \\ A \vdash{}}
  {\vdash}

  \and

  \inferrule*[right=$\mu$]
  {
    [A \vdash{}] \\\\ \vdots \\\\ \vdash
    %\inferrule*[fraction={~~~}]
    %{[A \vdash{}] \\\\ \vdots}
    %{\vdash}
  }
  {{}\vdash A}

  \and

  \inferrule*[right=$\tilde\mu$]
  {[{}\vdash A] \\\\ \vdots \\\\ \vdash}
  {A \vdash{}}
\end{mathpar}

The rules for logical connectives describe how to \emph{prove} and how to
\emph{refute} a formula whose principal connective is that connective.
These correspond strongly with the right and left rules, respectively, of LJ,
and for this reason, $\mu\tilde\mu$ is usually described elsewhere as a
sequent calculus.
For example, we could choose the following rules for disjunction.
To prove $A \vee B$, we can assume that both $A$ and $B$ are false, and derive
a contradiction.
To refute $A \vee B$, we can refute $A$ and $B$ separately.

\begin{mathpar}
  \inferrule*[right=$\vee$-r]
  {[A \vdash{}][B \vdash{}] \\\\ \vdots \\\\ \vdash}
  {{}\vdash A \vee B}

  \and

  \inferrule*[right=$\vee$-l]
  {A \vdash{} \\ B \vdash{}}
  {A \vee B \vdash{}}
\end{mathpar}

\subsection{Intro-elim sequent calculus: BBdPH}

Term assignment system introduced in~\cite{BBdPH93}.

  \section{Generic syntax}

\chapter{Linearity (the struggle)}
  \section{Intuitionistic linear logic}
  \section{Mechanisation survey}
  \section{DILL}
  DILL provides a syntax in which $\oc$ is characterised by a universal
  property.
  There is no need for any rule to depend on the entire context and rebind
  variables (as done by BBdPH), so is more ND-like.
  DILL makes $\oc$ behave more like a positive type former (see the deep
  identity function $\oc A \vdash \oc A$).
  DILL has been extended to dependent types: Vakar15.

  The idea of split contexts can also be applied to left-/right-rule systems.
  (Define this DILL variant.)

  \section{Notice what additive/multiplicative/exponential rules look like}

\chapter{Usage restriction via semirings}
  \section{Why semirings?}
  Addition gives the structural rules.
  Multiplication gives modalities.

  Compare with other work.
  About here, I can talk about Granule-style $\otimes$-pattern matching.

  \section{What are linear renaming and substitution?}
  \section{$\lambda\mathcal{R}$}

  \section{Linear natural deduction}
  This should stay fairly general, so it could be compiled to a DILL-style
  calculus as well as the one with usage annotations.
  Our framework compiles a ND system down to a sequent calculus.
  It isolates the building blocks of typing rules ($\vdash$, $\wedge$, $*$,
  $\Box$, $\cdot$) so that we can do more generic programming (e.g., usage
  checker).

\chapter{Semantics in worldly relations}

\chapter{Weighted multicategories}

\chapter{A framework for usage-restricted calculi}
  \def\prefix{../generic-lr/src/latex}

\CatchFileBetweenTags{\exItypes}%
{\prefix/Generic/Linear/Example/PaperExamples.tex}{exItypes}
\CatchFileBetweenTags{\exIlabels}%
{\prefix/Generic/Linear/Example/PaperExamples.tex}{exIlabels}
\CatchFileBetweenTags{\exIfunrules}%
{\prefix/Generic/Linear/Example/PaperExamples.tex}{exIfunrules}
\CatchFileBetweenTags{\exIsumrules}%
{\prefix/Generic/Linear/Example/PaperExamples.tex}{exIsumrules}

\CatchFileBetweenTags{\Premises}%
{\prefix/Generic/Linear/Syntax.tex}{Premises}
\CatchFileBetweenTags{\Rule}%
{\prefix/Generic/Linear/Syntax.tex}{Rule}
\CatchFileBetweenTags{\System}%
{\prefix/Generic/Linear/Syntax.tex}{System}

\CatchFileBetweenTags{\semp}%
{\prefix/Generic/Linear/Syntax/Interpretation.tex}{semp}
\CatchFileBetweenTags{\semr}%
{\prefix/Generic/Linear/Syntax/Interpretation.tex}{semr}
\CatchFileBetweenTags{\sems}%
{\prefix/Generic/Linear/Syntax/Interpretation.tex}{sems}

\CatchFileBetweenTags{\SimplePremises}%
{\prefix/Generic/Simple/Syntax.tex}{SimplePremises}

\section{Building up}

We assume that a type system comprises a set of (unparametrised) rules, each
of which has a conclusion and several premises containing subterms.
The primary investigation of this work is into what form the premises can take
while maintaining useful features of syntax.
We shall start with simple forms, allowing just for multiple subterms, and
then build on resource-sensitive bunches, variable binding, and modalities.

\System{}
\Rule{}

Given some way \AgdaFunction{⟦\AgdaUnderscore{}⟧p} of interpreting
\AgdaDatatype{Premises} in a \AgdaRecord{Ctx}, we can interpret a
\AgdaDatatype{Rule} against a conclusion and a context by checking that its
stated conclusion matches and then interpreting its premises.
Then, the entire \AgdaDatatype{System} can be interpreted by picking a rule
label (including parameters) \AgdaBound{l} and interpreting the selected rule
\AgdaBound{rs}\AgdaSpace{}\AgdaBound{l}.

\semr{}
\sems{}

\subsection{The language of Cartesian products}

\SimplePremises

\section{Bunched functions}

Suppose we have a left $R$-semimodule $M$.
Then we can define the following connectives on indexed type families.

\begin{align*}
  \top~z &:= \top \\
  (P \wedge Q)~z &:= P~z \wedge Q~z \\
  (P \to Q)~z &:= P~z \to Q~z \\
  I~z &:= z \subres 0 \\
  (P * Q)~z &:= \exists x,y.~z \subres x+y \wedge P~x \wedge Q~y \\
  (P \wand Q)~y &:= \forall x,z.~z \subres x+y \to P~x \to Q~z \\
  (r \cdot P)~z &:= \exists x.~z \subres rx \wedge P~x
\end{align*}

\paragraph{Agda definitions}
To follow.

The first six of these connectives are somewhat standard, the first three
corresponding to the usual sharing connectives and the second three
corresponding to the usual separating connectives.
The final connective, $r \cdot P$, is new.
We choose $\exists$ rather than $\forall$ because scaling will usually occur in
premises, i.e., to the left of arrows, and we want to avoid anything
higher-order.

\section{Generic syntax}

Rules $R$, premises $P$ and $Q$.

\begin{align*}
  R &::= P/A \\
  P,Q &::= \grR\Theta \vdash A \mid \top \mid P \wedge Q \mid I \mid P * Q
        \mid \gr r \cdot P
\end{align*}

Example rules:

\begin{itemize}
  \item With-introduction: ${\vdash A} \wedge {\vdash B} / A \with B$
  \item Annotated arrow introduction:
    ${\gr rA \vdash B} / \gr rA \multimap B$
  \item Annotated arrow elimination:
    ${\vdash \gr rA \multimap B} * \gr r \cdot {\vdash A} / B$
  \item Cases of annotated sum:
    ${\vdash \gr rA \oplus \gr sB}
    * ({\gr rA \vdash C} \wedge {\gr sB \vdash C}) / C$
\end{itemize}

\paragraph{Agda example}

We start by defining our types \AgdaDatatype{Ty}, which in this case will be the
conclusions of our judgements.
\AgdaFunction{Ann} is the type of usage annotations, which for now is an
arbitrary skew semiring.
We also have a type of rule labels \AgdaDatatype{`Sys}, which specifies all of
the parameters of a typing rule except for its conclusion and premises.
For example, in the \AgdaInductiveConstructor{`case} label, we need to fill in
the two summands (types \AgdaBound{A} and \AgdaBound{B}), as well as the return
type \AgdaBound{C}.
The actual form of the rule will be specified later.

\exItypes{} \exIlabels{}

We are then ready to build a system \AgdaFunction{Sys} with conclusions
\AgdaDatatype{Ty}, usage annotations \AgdaFunction{Ann}, and labels
\AgdaDatatype{`Sys}.
The rule corresponding to each label is given in the following.

The \AgdaInductiveConstructor{`lam} rule has one premise, which binds one
variable with usage annotation \AgdaBound{r} and type \AgdaBound{A}.
The \AgdaInductiveConstructor{`app} rule has two premises combined with
separating conjunction.
Furthermore, the second premise (corresponding to the argument of the function)
is subject to scaling by \AgdaBound{r}.
Intuitively, this description means that we must have enough of each variable to
use some in building a function, and then use some more \AgdaBound{r} times to
build enough arguments.

\exIfunrules{}

In the rules for sums, we use an alternative notation for rule descriptions,
with \AgdaFunction{{---}{---}} being an infix version of
\AgdaInductiveConstructor{rule}.
The introduction rules \AgdaInductiveConstructor{`inl} and
\AgdaInductiveConstructor{`inr} are straightforward --- each has one
premise which binds no new variables.
The \AgdaInductiveConstructor{`case} rule is more complicated.
The premises are first split into the eliminand and the continuations, using the
separating conjunction we saw for \AgdaInductiveConstructor{`app}.
However, the continuations are connected via the sharing conjunction, reflecting
the fact that in any given run of the program, only one of the branches will be
taken, so usages from each individually should be added to usages of the
eliminand.

\exIsumrules{}

\paragraph{Interpretation as syntax}
We can say something like ``premise connectives are interpreted as the
corresponding bunched connectives, where appropriate''.

\begin{align*}
  \sem{P/A} &:= \sem P \to {- \vdash A} \\
  \sem{\grR\Theta \vdash A} &:= {-, \grR\Theta \vdash A} \\
  \sem{\top} &:= \top \\
  \sem{P \wedge Q} &:= \sem P \wedge \sem Q \\
  \sem{I} &:= I \\
  \sem{P * Q} &:= \sem P * \sem Q \\
  &\ldots
\end{align*}

\paragraph{Agda version}
The interpretation of a system is the selection of a rule, together with the
interpretation of that rule.

\sems{}

The interpretation of a rule is a check that the rule targets the desired
conclusion, together with the interpretation of the rules premises.

\semr{}

For the premises, connectives are interpreted in the obvious way.
Premises can ask for subterms via the $\AgdaInductiveConstructor{\_`⊢\_}$
constructor, which supplies the new variables (in particular, each variable's
type and usage annotation) \AgdaBound{Γ} and the desired conclusion
\AgdaBound{A} of the subterm.

\semp{}

\section{Values and computations}

Following \cite{AACMM20}, we call any semantic realisation of a variable a
\emph{value}, and any semantic realisation of a term a \emph{computation}.
Scoped families of values are named $\mathcal V$, while scoped families of
computations are named $\mathcal C$.
\todo{This wording is difficult.}
Values and computations are concepts with attitudes.

\section{Environments}

An environment is a linear map from variables in one context to values in
another.

\missingfigure{Agda definition of \AgdaRecord{\_{---}Env}}

An environment in \AgdaBound{PΓ} for \AgdaBound{$\mathcal V$}-values in
\AgdaBound{QΔ} is a linear map \AgdaField{M} such that $Q \le PM$ and for any
$P'$ and $Q'$ such that $Q' \le P'M$, a function from variables in $P'Γ$ to
\AgdaBound{$\mathcal V$}-values in $Q'Δ$.

\section{Thinnings}

\todo[inline]{I think thinnings should be called \emph{renamings}.}

In a usage-annotated sequent calculus, their are 4 basic structural rules:
exchange, weakening of $0$-annotated formulae, contraction of $+$-annotated
formulae, and subusaging.
Their typical forms are given below.

\begin{mathpar}
  \inferrule*[right=exch]
  {\grP\Gamma, \gr sB, \gr rA, \grQ\Delta \vdash Z}
  {\grP\Gamma, \gr rA, \gr sB, \grQ\Delta \vdash Z}
  \and
  \inferrule*[right=weak]
  {\grP\Gamma, \grQ\Delta \vdash Z}
  {\grP\Gamma, 0A, \grQ\Delta \vdash Z}
  \and
  \inferrule*[right=cont]
  {\grP\Gamma, \gr rA, \gr sA, \grQ\Delta \vdash Z}
  {\grP\Gamma, (\gr r + \gr s)A, \grQ\Delta \vdash Z}
  \and
  \inferrule*[right=subuse]
  {\grP \le \grQ \\\\ \grQ\Gamma \vdash Z}
  {\grP\Gamma \vdash Z}
\end{mathpar}

Generalising away from $\vdash$, it can be helpful to know when an arbitrary
scoped family respects these structural rules.
First, we note that the structural rules of sequent calculus correspond
directly to \emph{renaming} in natural deduction.
A renaming is an environment of variables, i.e., a \emph{linear} map from
variables to variables.
The linearity is what gives us the restrictions on annotations in the weakening
and contraction rules.

\missingfigure{Example renaming}

%A thinning is an environment of syntactic variables.
%In effect, thinning allows for permutation and slackening of existing
%variables, and introduction of new discardable variables at any position.
%These actions are in direct correspondence with the usual sequent calculus
%structural rules of exchange, contraction, and weakening, as well as subusaging.

Given this definition, $\square$ and \texttt{Thinnable} are as in AACMM.
\texttt{Kripke} is modified to use separating implication.

\section{Substitution}

\todo[inline]{This seems like a good place to introduce substitution, straight
after renaming, though we don't even have terms yet.}
Whereas a renaming is a linear map from variables to variables, a substitution
is a linear map from variables to terms.

\section{A layer of syntax is functorial}


\chapter{Generic programs}
  \section{Usage checker}

\chapter{Investigations that are now easier}
  \section{Linear/non-linear logic}
  We can express Benton's linear/non-linear logic~\cite{Benton94} in the
framework.
There are two apparent discrepancies between L/nL and what we have seen so far
to be expressible in the framework: the existence of multiple judgement modes
and the restriction on the kinds of assumptions based on the mode.
We will start by addressing these discrepancies, before presenting the encoding
in full and proving the resulting system logically equivalent to
$\lambda\gr{\mathcal R}$.

\subsection{Encoding L/nL}

In L/nL, we have a \emph{linear} and an \emph{intuitionistic} mode, with
separate types (respectively $A$ and $X$) and typing judgements (respectively
$\vdash_{\mathcal L}$ and $\vdash_{\mathcal C}$) connected by modalities $F$ and
$G$.
The mode of the conclusion type is the same as the mode of the judgement, so
only a distinction in the conclusion mode is necessary.
To achieve this distinction in the types, I have an indexed type family
$\mathrm{Ty} : \mathrm{Frag} \to \mathrm{Set}$, and set the framework types to
be the type $(f : \mathrm{Frag}) \times \mathrm{Ty}\,f$.
% \AgdaDatatype{Ty}\AgdaSymbol{:}\AgdaDatatype{Frag}\AgdaSymbol{$\to$}
% \AgdaDatatype{Set}
I present these judgements in a standard paper style in \cref{fig:lnl-types}.

In L/nL, while linear judgements can have both linear and intuitionistic
assumptions, intuitionistic judgements can only have intuitionistic assumptions.
I encode this invariant on intuitionistic judgements using the $\Box$
premise combinator on all intuitionistic conclusions and wherever
intuitionistic premises occur for linear conclusions, as seen in
\cref{fig:lnl-desc}.
This scheme means that the rules contain the maximal number of boxes, which is
convenient when consuming terms.
For ease of \emph{producing} terms, we would want a ``garbage in; garbage out''
style with a minimal number of boxes, which could be achieved in two different
ways:
\begin{itemize}
  \item Ensure that $\Gamma, \gr1\,\lin A \vdash \intu X$ is never derivable.
    This leaves boxes only on rules $1$-I and $G$-I.
  \item Ensure that, when working bottom-up, we never transition from a linear
    to an intuitionistic conclusion with linear assumptions present.
    This leaves boxes only on rules $F$-I and $G$-E.
\end{itemize}
It would be interesting to show the equivalence of these three approaches, but
that is left to future work.

\begin{figure}
  \begin{align*}
    A, B, C &\Coloneqq I \mid A \otimes B \mid A \multimap B \mid FX \\
    X, Y, Z &\Coloneqq 1 \mid X \times Y \mid X \to Y \mid GA \\
    \mathcal J &\Coloneqq \lin A \mid \intu X
  \end{align*}
  \caption{Types and judgement forms of L/nL}
  \label{fig:lnl-types}
\end{figure}

\begin{figure}
  \begin{mathpar}
    % lin rules
    \ebrule[comb]{%
      \hypo{I^*}
      \infer1[$I$-I]{\vdash \lin I}
    }
    \and
    \ebrule[comb]{%
      \hypo{\vdash \lin I}
      \hypo{*}
      \hypo{\vdash \lin C}
      \infer3[$I$-E]{\vdash \lin C}
    }
    \and
    \ebrule[comb]{%
      \hypo{\vdash \lin A}
      \hypo{*}
      \hypo{\vdash \lin B}
      \infer3[$\otimes$-I]{\vdash \lin A \otimes B}
    }
    \and
    \ebrule[comb]{%
      \hypo{\vdash \lin A \otimes B}
      \hypo{*}
      \hypo{\gr1\,\lin A, \gr1\,\lin B \vdash \lin C}
      \infer3[$\otimes$-E]{\vdash \lin C}
    }
    \and
    \ebrule[comb]{%
      \hypo{\gr1\,\lin A \vdash \lin B}
      \infer1[$\multimap$-I]{\vdash \lin A \multimap B}
    }
    \and
    \ebrule[comb]{%
      \hypo{\vdash \lin A \multimap B}
      \hypo{*}
      \hypo{\vdash \lin A}
      \infer3[$\multimap$-E]{\vdash \lin B}
    }
    \and
    \ebrule[comb]{%
      \hypo{\Box\plr{\vdash \intu X}}
      \infer1[$F$-I]{\vdash \lin FX}
    }
    \and
    \ebrule[comb]{%
      \hypo{\vdash \lin FX}
      \hypo{*}
      \hypo{\gr\omega\,\intu X \vdash \lin C}
      \infer3[$F$-E]{\vdash \lin C}
    }
    \and
    % int rules
    \ebrule[comb]{%
      \hypo{\Box\plr{\dot1}}
      \infer1[$1$-I]{\vdash \intu 1}
    }
    \and
    \text{(no $1$-E)}
    \and
    \ebrule[comb]{%
      \hypo{\Box(\vdash \intu X}
      \hypo{\dottimes}
      \hypo{\vdash \intu Y)}
      \infer3[$\times$-I]{\vdash \intu X \times Y}
    }
    \and
    \ebrule[comb]{%
      \hypo{\Box\plr{\vdash \intu X \times Y}}
      \infer1[$\times$-El]{\vdash \intu X}
    }
    \and
    \ebrule[comb]{%
      \hypo{\Box\plr{\vdash \intu X \times Y}}
      \infer1[$\times$-Er]{\vdash \intu Y}
    }
    \and
    \ebrule[comb]{%
      \hypo{\Box\plr{\gr\omega\,\intu A \vdash \intu B}}
      \infer1[$\to$-I]{\vdash \intu A \to B}
    }
    \and
    \ebrule[comb]{%
      \hypo{\Box(\vdash \intu A \to B}
      \hypo{\dottimes}
      \hypo{\vdash \intu A)}
      \infer3[$\to$-E]{\vdash \intu B}
    }
    \and
    \ebrule[comb]{%
      \hypo{\Box\plr{\vdash \lin A}}
      \infer1[$G$-I]{\vdash \intu GA}
    }
    \and
    \ebrule[comb]{%
      \hypo{\Box\plr{\vdash \intu GA}}
      \infer1[$G$-E]{\vdash \lin A}
    }
  \end{mathpar}
  \caption{Description of the logical rules of L/nL}
  \label{fig:lnl-desc}
\end{figure}

\begin{proposition}
  We can construct the following translations.
  \begin{align}
    (\Theta \vdash_{\mathcal C} X) &\to (\gr\omega\Theta \vdash \intu X) \\
    (\Theta; \Gamma \vdash_{\mathcal L} A) &\to
    (\gr\omega\Theta, \gr1\Gamma \vdash \lin A)
  \end{align}
\end{proposition}

\begin{proposition}
  We can construct the following translations, where $\Gamma_{\gr r}$ is the
  list of types in $\Gamma$ that are annotated $\gr r$.%
  \todo{I'm not sure the resultant contexts are well formed.}
  \begin{align}
    (\Gamma \vdash \intu X) &\to (\Gamma_{\gr\omega} \vdash_{\mathcal C} X) \\
    (\Gamma \vdash \lin A) &\to
    (\Gamma_{\gr\omega}; \Gamma_{\gr1} \vdash_{\mathcal L} A)
  \end{align}
\end{proposition}

\subsection{Translating between L/nL and $\lambda\gr{\mathcal R}$}

The translations largely follow Benton's originals.
In this section, I take $\name$ to be the system described in
\cref{fig:lr-bunched} instantiated to the $\{\gr0 > \gr\omega < \gr1\}$
posemiring and restricted to the fragment with only connectives $I$, $\otimes$,
$\multimap$, and $\oc\gr\omega$.
Notably, this excludes $\oc\gr0$ and $\oc\gr1$.
I write $\oc\gr\omega$ as just $\oc$, as in traditional linear logic.
Under these restrictions, Benton's translations of types can be used verbatim
(\cref{fig:lnl-lr-types}).

\begin{figure}
  \centering
  \begin{subfigure}{.49\linewidth}
    \centering
    \begin{align*}
      (-)^\circ : \mathrm{Ty}_{\name} \to \mathrm{Ty}_{\lin} \\
      \begin{aligned}
        I^\circ &= I \\
        \plr{A \otimes B}^\circ &= A^\circ \otimes B^\circ \\
        \plr{A \multimap B}^\circ &= A^\circ \multimap B^\circ \\
        \plr{\oc A}^\circ &= GFA^\circ
      \end{aligned}
    \end{align*}
  \end{subfigure}
  \begin{subfigure}{.49\linewidth}
    \centering
    \begin{align*}
      (-)^* : \mathrm{Ty}_\lnl \to \mathrm{Ty}_{\name} \\
      \begin{aligned}
        I^* &= I \\
        \plr{A \otimes B}^* &= A^* \otimes B^* \\
        \plr{A \multimap B}^* &= A^* \multimap B^* \\
        \plr{FX}^* &= \oc X^* \\
        1^* &= I \\
        \plr{X \times Y}^* &= \oc X^* \otimes \oc Y^* \\
        \plr{X \to Y}^* &= \oc X^* \multimap Y^* \\
        \plr{GA}^* &= A^*
      \end{aligned}
    \end{align*}
  \end{subfigure}
  \caption{Translation of types between L/nL and $\name$}
  \label{fig:lnl-lr-types}
\end{figure}

We extend both $(-)^\circ$ and $(-)^*$ to contexts pointwise on the types.
For $(-)^\circ$, this means that the translation lands outside what is usually
expressible in L/nL whenever the context contains annotation $\gr\omega$.
For example, $\plr{\gr\omega I}^\circ = \plr{\gr\omega\,\lin I}$, and we do not
see ``$\gr\omega\,\lin$'' in translations from Benton's L/nL.
This could be fixed by inserting $G$ in such cases.
As for $(-)^*$, note that we do not insert a $\oc$ on intuitionistic assumptions
as Benton does.
Instead, the annotation $\gr\omega$ achieves the same effect for us.
This discrepancy essentially amounts to the fact that $\name$ is more like DILL
than Girard's ILL, and Benton's translations act upon the latter.

\begin{theorem}
  We can translate from $\name$ to the linear fragment of L/nL.
  \begin{align}
    (\Gamma \vdash_{\name} A) \to (\Gamma^\circ \vdash_\lnl \lin A^\circ)
  \end{align}
\end{theorem}

\begin{theorem}
  We can translate any L/nL term to a $\name$ term as follows.
  \begin{align}
    (\Gamma \vdash_\lnl \intu X) &\to (\Gamma^* \vdash_{\name} X^*) \\
    (\Gamma \vdash_\lnl \lin A) &\to (\Gamma^* \vdash_{\name} A^*)
  \end{align}
\end{theorem}

  \section{$\mu\tilde\mu$-calculus}
  \section{Graphs}

\chapter{Conclusion}



%%%%%%%%%%%%%%%%%%%%%%%%%%%%%%%%%%%%%%%%%%%%%%%%%%%%%%%%%%%%%%
\appendix
\chapter{Stuff That Didn't Fit Anywhere Else}
%%%%%%%%%%%%%%%%%%%%%%%%%%%%%%%%%%%%%%%%%%%%%%%%%%%%%%%%%%%%%%


%%%%%%%%%%%%%%%%%%%%%%%%%%%%%%%%%%%%%%%%%%%%%%%%%%%%%%%%%%%%%%
\addcontentsline{toc}{chapter}{Bibliography}
\bibliographystyle{alpha}
\bibliography{quantitative}

\end{document}
